\documentclass[10pt]{article}
\usepackage[margin=20mm]{geometry}
\usepackage{amsmath,amssymb,mathtools,fontspec,unicode-math,xurl,hyperref}
\IfFontExistsTF{DejaVu Serif}{\setmainfont{DejaVu Serif}}{\setmainfont{Latin Modern Roman}}
\setmathfont{latinmodern-math.otf}
\setlength{\parindent}{0pt}\setlength{\parskip}{6pt}\sloppy
\title{A Technical Review of
Central Limit Theorem for Rényi Divergence of Infinite Order}
\author{AI-assisted technical review}
\date{}
\begin{document}
\maketitle
Abstract: The Esscher-transform criterion for convergence in infinite-order Rényi divergence is checked in the acquired arXiv version. The critical-zone argument, its periodic specialization, and a reduction covering an arbitrary initial smoothing index are supported by the calculations below. Several displayed identities have incorrect signs or normalizing factors. Explicit Gaussian and normalization tests establish these local errors; the intended corrections preserve the checked proof. No false final theorem is demonstrated. The assessment is Minor, with the cited external local-limit estimates remaining dependencies rather than independently certified results.\par
\section*{Identity, hypotheses, and coverage}
The source is the 29-page arXiv:2402.02259v2 by Sergey G. Bobkov and Friedrich Götze, dated 19 June 2024 [paper]. Its associated journal DOI is 10.1214/24-AOP1711. This report applies to that acquired preprint; equivalence with the journal-final text has not been established. Page numbers below are one-based PDF pages and agree with the printed numbers.\par
All substantive source text, Sections 1–11 on pp. 1–27, was read. The references on pp. 27–29 were read as bibliographic entries, not as their underlying papers. The main proof, Esscher identities, explicit periodic examples, and the weighted-convolution argument were checked. The quantitative Fourier-tail estimate imported in Lemma 6.1 and the refined Richter theorem quoted as Theorem 10.2 were not independently re-proved from their original publications. Bounds that use those results retain that dependency. No author code was executed. Bounded source and public-metadata searches identified no separate computational supplement or correction notice; this does not prove their absence. The official arXiv record describes v2 as correcting typos and updating Section 11 [metadata]; it does not certify every displayed formula in v2.\par
Write \(Z_n=n^{-1/2}\sum_{j=1}^nX_j\), where the \(X_j\) are independent copies of a real random variable with mean zero and variance one. The operative assumptions before Theorem 2.1, p. 4, are \[p_{n_0}\le C\varphi\quad\hbox{a.e. for some finite }n_0,C,
 \qquad L(t)=\mathbb E e^{tX}\le e^{t^2/2}\quad(t\in\mathbb R).\] The first is a Gaussian envelope after finitely many summands, stronger than existence of a density. Put \(K=\log L\), \(A=t^2/2-K\ge0\), and \(T_\infty(p\Vert\varphi)=\mathop{\rm ess\,sup}(p/\varphi)-1\). The claimed criterion is \[\lim_{A(t)\to0}\max(A''(t),0)=0.                         \tag{1.1}\] It includes the values at every zero of \(A\), and sequences approaching infinity with \(A(t)\to0\). It is not only a condition near the origin. The Gaussian case \(A\equiv0\) satisfies it. The smallest-positive-period assumption in Theorem 2.3 excludes a constant periodic factor; a Gaussian should instead be handled directly.\par
\section*{Checks of the main argument}
\subsection*{Tilt, moments, and the critical zone}
For \(Q_hp(x)=e^{hx}p(x)/L(h)\), direct integration gives \[L_{Q_hp}(t)=\frac{L(t+h)}{L(h)},\quad
 m_h=K'(h),\quad
 \sigma_h^2=K''(h)=\frac{L(h)L''(h)-L'(h)^2}{L(h)^2}>0.       \tag{2.1}\] Positivity uses nondegeneracy of the tilted distribution. Its support is the same as that of \(X\). Convolution and scaling of this probability tilt give, with \(x,h\in\mathbb R\), \[\frac{p_n(x\sqrt n)}{\varphi(x\sqrt n)}
 =\sqrt{2\pi}\exp\{n(x-h)^2/2-nA(h)\}
     Q_{h\sqrt n}p_n(x\sqrt n).                            \tag{2.2}\] This independently reproduces (3.7), p. 7.\par
Since \(A''\le1\) and \(A\ge0\), the quadratic upper Taylor bound, minimized over its increment, gives \((A')^2\le2A\). The elementary density–variance bound \(M^2\operatorname{Var}\ge1/12\), combined with the Gaussian envelope, gives \[\sigma_h^{-2}\le(6/\pi)C^2e^{2A(h)}.                       \tag{2.3}\] Consequently the tilted variance stays uniformly away from zero in the critical zones. The tilted third-moment bound of Lemma 5.2 and the stated local limit estimate lead to \[\frac{p_n(x\sqrt n)}{\varphi(x\sqrt n)}
 =\frac{\exp\{-nA(x)-nv_x^2/2\}}{\sigma_x}
       +O(C^4n^{-1/2}),\qquad v_x=\frac{A'(x)}{\sigma_x},   \tag{2.4}\] uniformly when \(A(x)\le a/(n-1)\) and \(n\ge4(a+1)\), as in Lemma 7.1, pp. 14–15. Outside that set, the envelope in (4.3), p. 9, decays as \(C\sqrt2e^{-(n-1)A(x)}\).\par
These bounds justify sufficiency in Section 8: first fix the critical-zone parameter, let \(n\) tend to infinity, and then enlarge that parameter. For necessity, (2.3) and \((A')^2\le2A\) bound \(v_x^2\) by a constant times \(C^2A(x)\). Applying (2.4) on arbitrarily small critical zones forces \(\sigma_x^{-2}\le1+o(1)\), hence (1.1). In the first display on p. 17 the exponential inherited from (2.3) should be \(e^{2A(x)}\). For \(a\le1\) and \(n\ge3\), \(A(x)\le1/2\) and \(12e/\pi<12\); the following bound \(v_x^2\le12C^2A(x)\) is therefore recovered. This correction does not alter the asymptotic argument.\par
Strict subgaussianity itself is necessary: \(p_n\le(1+o(1))\varphi\) implies \(L(h)^n\le(1+o(1))e^{nh^2/2}\) for every fixed \(h\), by evaluating the moment-generating function at \(h\sqrt n\). Taking logarithms and dividing by \(n\) gives \(A(h)\ge0\).\par
\subsection*{The initial smoothing index and all integers \(n\)}
Section 8 assumes \(n_0=1\) for simplicity. The following supplies the reduction for its stated general hypothesis. Let \(m=n_0\). The density of \(S_m=\sum_{j=1}^mX_j\) is bounded by \(C\) times the \(N(0,m)\) density. Its tilt has maximum at most \(C(2\pi m)^{-1/2}e^{mA(h)}\). Convolving with the tilted probability law of \(X\) cannot increase that maximum. Using (2.2) for \(N=m+1\) therefore gives \[\frac{p_N(h\sqrt N)}{\varphi(h\sqrt N)}
 \le C\sqrt{N/m}\,e^{-A(h)}\le C\sqrt{N/m}.                \tag{2.5}\] Thus both \(Z_m\) and \(Z_{m+1}\) have finite \(T_\infty\). Their deficit functions are \(A_r(t)=rA(t/\sqrt r)\), with \(A_r''(t)=A''(t/\sqrt r)\). Condition (1.1) is preserved.\par
Apply the verified \(n_0=1\) argument to both block laws. Every sufficiently large \(N\) can be written \(N=am+b(m+1)\), where \(b\equiv N\pmod m\) is chosen near \(N/[2(m+1)]\); then both component lengths are of order \(N\). Their normalized densities are bounded by \((1+\epsilon_N)\varphi\) and \((1+\delta_N)\varphi\), with both errors tending to zero. Independence and convolution of the corresponding Gaussian bounds yield \(p_N\le(1+\epsilon_N)(1+\delta_N)\varphi\). This proves convergence along every integer, not only multiples of \(m\). The same balanced decomposition preserves a bound \(O((\log N)^3/N)\) when that bound has been established for both block laws; the separation and periodic hypotheses are also preserved by the block transformation.\par
\subsection*{An explicit valid periodic example}
For Example 9.5, p. 20, put \(P(t)=(1-4\sin^2t)^2\sin^4t\). Direct expansion gives \[P(t)=\frac94-\frac{15}4\cos2t+\frac{17}8\cos4t
                   -\frac34\cos6t+\frac18\cos8t.\] Let these coefficients be \(a_k\), including \(k=0\), and set \[B=\sum_{k\in\{0,2,4,6,8\}}|a_k|e^{k^2/2},\quad c=(2B)^{-1},
 \qquad q(x)=1-c\sum_k a_ke^{k^2/2}\cos(kx).\] Then \(1/2\le q\le3/2\), and \(p=q\varphi\) is a probability density with \(L(t)=e^{t^2/2}(1-cP(t))\). This follows by integrating each Gaussian cosine term. Because \(P(0)=P'(0)=P''(0)=0\), the mean and variance are zero and one, respectively. Since \(P\ge0\), strict subgaussianity holds, and \(T_\infty(p\Vert\varphi)\) is finite.\par
Writing \(P=Q^2\), \(Q=\sin^2t-4\sin^4t\), gives \(Q(\pi/6)=0\), \(Q'(\pi/6)=-\sqrt3/2\), and \(\Psi''(\pi/6)=-3c/2\) for \(\Psi=1-cP\). At this contact point, (2.4) yields the limiting density ratio \[(1-3c/2)^{-1/2}>1.\] Thus the source’s failure-of-convergence example is valid after the displayed normalization correction below. The contrasting \(\sin^4t\) example has only the endpoint contacts in a fundamental period and vanishing second derivative there.\par
\section*{A demonstrated local problem and its repair}
Problem 1 (Displayed signs and normalizations).  Several formulas in the acquired v2 are false literally. Their corrections are determined by the definitions and can be tested without any unverified example.\par
For the Gaussian tests below take \(X\sim N(0,1)\), so all mean, variance, strict-subgaussianity, density and finite-divergence assumptions hold. Set \(h=1\) when a nonzero tilt is required.\par
1. On p. 6, after Definition 3.1, the identity \(Q_h\varphi(x)=\varphi(x+h)\) must read \(\varphi(x-h)\). Completing the square gives a tilted mean \(h\), whereas the displayed right side has mean \(-h\).\par
2. On p. 7, the proof of Proposition 3.2 divides \(L(t+h)\) by \(L(t)\); every normalizing denominator in that display must instead be evaluated at \(h\). At \(t=0,h=1\), the printed Gaussian expression is \(e^{1/2}\), whereas a probability moment-generating function at zero equals one. Equation (5.1), p. 10, already has the correct denominator.\par
3. On p. 10, the displayed variance numerator omits the factor \(L(h)\) multiplying \(L''(h)\). For the Gaussian at \(h=1\), the printed fraction equals \(2e^{-1/2}-1\), instead of the true variance one. Formula (2.1) gives the correction.\par
4. On p. 3, the periodic Laplace factor must be \(e^{+t^2/2}\), as correctly used in (9.1) and (9.4). On p. 13, in the proof of Lemma 6.1, the normal characteristic-function factor must instead be \(e^{-t^2/2}\) in all three displays. With the printed positive exponent, its difference from a bounded characteristic function is not integrable over the real line.\par
5. On p. 20, Example 9.5 says \(L(t)=1-c\Psi(t)\). Its value at zero is \(1-c\), violating normalization for the stated \(c>0\). The correct formula is \(L(t)=e^{t^2/2}\Psi(t)\), consistent with p. 19 and the explicit positive-density construction above.\par
6. On p. 25, the \(c_0=0\) convolution has \(q_2(y/c_1)\) and then \(\varphi(y/c_1)\). Both denominators must be \(c_2\). With independent standard Gaussians, \(c_1=\sqrt3/2,c_2=1/2\), the printed expression has integral \(\sqrt3\); the corrected convolution has integral one and variance \(c_1^2+c_2^2=1\).\par
\paragraph*{Counterevidence, dependency, and severity.}\mbox{}\par
Definition (3.1), formula (3.7), the correct tilt formula (5.1), the use of \(K''\) in Section 8, and the consistent periodic definition (9.1) identify the intended expressions. Direct substitution of the corrected formulas restores the checked algebra. These are genuine false displayed identities, but they do not establish failure of the central limit criterion or the periodic examples. The assessment is Minor. The repair is verified locally; author adoption and any changes in the journal version are unknown.\par
Two additional scope details do not warrant a false-theorem claim. The critical zone in (4.5), p. 10, should range over all real \(h\), as it explicitly does in the actual proof on p. 16. Also, signed coefficients in Corollary 11.2 need not be assumed positive without explanation: absorb each sign into its summand, retaining independence, strict subgaussianity and the same Gaussian density bound. Identical distribution is unnecessary in the application of Lemma 11.3. Its contour shift, uniform Riemann–Lebesgue argument and weighted Plancherel bounds on pp. 25–27 then give the stated envelope.\par
\section*{Assessment and limits}
No Major defect or false final theorem was established. The main criterion and the explicit positive and negative periodic tests pass the checks recorded here, with the local corrections in Problem Problem 1 (\#prob:typos). The refined rate statements remain conditional on the externally quoted Richter estimate, whose hypotheses are matched in the source but whose original proof was not audited. One-sided relative error is essential: for any compactly supported density, the absolute relative error equals one outside the support of each normalized sum.\par
Scientific adjudication remains pending and no human reviews have been supplied. This completed initial report is not independent acceptance. Its source, exact TeX and compiled PDF, actual page-inspection receipts, symbolic calculation and outputs, compiler/dependency payloads and current support inventory are retained in `verification/0492`. The source-PDF SHA256 is 6a1a96abbb1061bdd8bec26b77fb3626c394481caa6c11d9e2c478bf7535c6cf.\par
9 S. G. Bobkov and F. Götze, Central Limit Theorem for Rényi Divergence of Infinite Order, arXiv:2402.02259v2, 19 June 2024. https://arxiv.org/abs/2402.02259v2. arXiv, version history and author comments for arXiv:2402.02259, observed 29 October 2026. https://arxiv.org/abs/2402.02259.\par
\end{document}
